\section{Tasks are too small}
\label{section:node-level:tasks-too-small}


If tasks (or parallel sections) are too small, we pay an excessive overhead.
The concurrency level might be brilliant, but that does not help us if we have too much task creation and task clean-up overhead.

\begin{hypothesis}
 There are enough tasks and they are not excessively synchronising each other, i.e.~there are not that many dependencies, but they are just too small and hence create too much administration overhead.
\end{hypothesis}


There is not really a definition of a ``too small task''.
However, we can use rule of thumbs from parallel loop segments and consult popular developer platforms and read about the recommendation that a reasonable task should at least run for tens of microseconds.
The extreme case of tiny tasks are empty tasks which some codes use to model dependencies, i.e.~to synchronise sets of tasks in a task graph.
If we have too many of those guys, the fine granular decomposition of a problem into a task graph becomes a performance problem in its own right.


\subsection{Data collection}

Different to excessive synchronisation (Section~\ref{section:node-level:synchronisation}), spotting too small tasks is often challenging.
Either, we do not see locks or excessive spinning at all, or we see these but within the task runtime itself.
After all, small tasks lead to a low overhead from a user code's point of view, as the overhead is typically attributed to the task runtime.

A plain profile is difficult to assess as well, as very short task functions do not appear high up on the list either.
They might have a very high call number, but the combination low runtime plus high call number also implies that they are often hard to distinguish from sole helper functions.

Due to the fact that this flaw is not trivial to spot, it does not come as a surprise that many task-based codes bring along their own task tracing.
The present manuscript cannot and should not cover these bespoke solutions, but rather advocates for a symptomatic analysis or, if possible, measurements through a runtime's tool interface.


\subsubsection{Symptom analysis with \openmp}
For systems employing \openmp, very small tasks materialise in a lot of time spent in \texttt{\_\_kmp\_fork\_barrier}.
When we open up the trace, we see that this function is called by \texttt{\_\_kmp\_launch\_thread}.
So all the overhead arises from \texttt{libiomp5.so} or similar.



\begin{instructions}{\vtune}
 In \vtune's hotspot analysis, you see in the summary a very high overhead time.
 Indeed, studies of the timeline confirm that we see that all the time is burn in spin and overhead. 
 
 In \vtune's Threading analysis as well as the HPC Performance Charaterization, we are informed on the summary page that the Effective CPU (or physical core) Utilization is low.
 Low here implies less than 80\%.
 Unfortunately, all the tools claim that the \openmp~region with the biggest potential gain is something similar to
 \texttt{unknown\$omp\$parallel}, and if we zoom into the area of interest, the threading analysis quickly tells use that there is no data to display.
\end{instructions}



\begin{wip}
 I haven't yet found a way to find out which tasks are too cheap. 
\end{wip}

\subsection{Evaluation}


\begin{enumerate}[noitemsep]
  \item[$\boxtimes$] There are too many tiny tasks modelled as real tasks.
    \begin{itemize}[noitemsep]
      \item[$\Rightarrow $] Inline tasks, i.e.~call the tasks as functions rather than spawning them as real tasks.
      \item[$\Rightarrow $] Randomise task creation and therefore throttle the task production.
    \end{itemize}
\end{enumerate}


\noindent
In \openmp, ensuring that logical tasks are not mapped onto physical tasks is very simple.
You can add a \texttt{if (xxx)} annotation to the pragma that switches it on or off depending on a boolean predicate.


% \paragraph*{How it works}
% 
% \begin{itemize}[noitemsep]
%   \item[$\boxplus$] Validate that the administrative (non-science) fraction of the code scales with the scientific core calculations.
% \end{itemize}
% 
% 
% \paragraph*{How it does not work}
% 
% \begin{itemize}[noitemsep]
%   \item[$\boxminus$] The administrative overhead will always increase relatively, so an increase in itself is not a flaw. An extraordinary increase coupled with limited scalability however is.
% \end{itemize}


